% from ega2-1-fr.tex
\oldpage[II]{33}

en composant l'homomorphisme
$M_{(f)}\otimes_{S_{(f)}}N_{(f)}\to M_f\otimes_{S_f}N_f$
(provenant des injections $M_{(f)}\to M_f$, $N_{(f)}\to N_f$ et
$S_{(f)}\to S_f$) et l'isomorphisme canonique
$M_f\otimes_{S_f}N_f\xrightarrow{\sim}(M\otimes_S N)_f$
(\hyperref[0.1.3.4-fr]{0, 1.3.4}) et en notant que, par définition du
produit tensoriel de deux modules gradués, ce dernier isomorphisme conserve
les degrés~; pour $x\in M_{md}$, $y\in N_{nd}$ ($m\geq0$, $n\geq0$), on a
donc
\[
  \lambda_f((x/f^m)\otimes(y/f^n))=(x\otimes y)/f^{m+n}.
\]

Il résulte aussitôt de cette définition que si $g\in S_e$ ($e>0$), le
diagramme
\[
  \xymatrix{
    M_{(f)}\otimes_{S_{(f)}}N_{(f)}\ar[r]^{\lambda_f}\ar[d]
    & (M\otimes_S N)_{(f)}\ar[d]\\
    M_{(fg)}\otimes_{S_{(fg)}}N_{(fg)}\ar[r]_{\lambda_{fg}}
    & (M\otimes_S N)_{(fg)}
  }
\]
(où la flèche verticale de droite est l'homomorphisme canonique et celle de
gauche provient des homomorphismes canoniques) est commutatif. On déduit
donc des $\lambda$ un homomorphisme canonique fonctoriel de
$\mathcal{O}_X$-Modules
\[
\label{II.2.5.11.2-fr}
  \lambda:\widetilde{M}\otimes_{\mathcal{O}_X}\widetilde{N}
  \longrightarrow\widetilde{M\otimes_S N}.
\tag{2.5.11.2}
\]

Considérons en particulier deux idéaux gradués $\mathfrak{J}$,
$\mathfrak{K}$ de $S$~; comme $\widetilde{\mathfrak{J}}$ et
$\widetilde{\mathfrak{K}}$ sont des faisceaux d'idéaux de
$\mathcal{O}_X$, on a un homomorphisme canonique
$\widetilde{\mathfrak{J}}\otimes_{\mathcal{O}_X}
\widetilde{\mathfrak{K}}\to\mathcal{O}_X$~; le diagramme
\[
\label{II.2.5.11.3-fr}
  \xymatrix{
    \widetilde{\mathfrak{J}}\otimes_{\mathcal{O}_X}
    \widetilde{\mathfrak{K}}\ar[rr]^\lambda\ar[dr]
    && \widetilde{\mathfrak{J}\otimes_S\mathfrak{K}}\ar[dl]\\
    & \mathcal{O}_X &
  }
\tag{2.5.11.3}
\]
est alors commutatif. On se ramène en effet à le vérifier dans chaque
ouvert $D_+(f)$ ($f$ homogène dans $S_+$) et cela résulte aussitôt de la
définition (\hyperref[II.2.5.11.1-fr]{2.5.11.1}) de $\lambda_f$ et de
(\hyperref[I.1.3.13-fr]{I, 1.3.13}).

Notons enfin que si $M$, $N$, $P$ sont trois $S$-modules gradués, le
diagramme
\[
\label{II.2.5.11.4-fr}
  \xymatrix{
    \widetilde{M}\otimes_{\mathcal{O}_X}\widetilde{N}
    \otimes_{\mathcal{O}_X}\widetilde{P}
    \ar[r]^{\lambda\otimes1}\ar[d]_{1\otimes\lambda}
    & \widetilde{M\otimes_S N}\otimes_{\mathcal{O}_X}\widetilde{P}
    \ar[d]^\lambda\\
    \widetilde{M}\otimes_{\mathcal{O}_X}\widetilde{N\otimes_S P}
    \ar[r]_\lambda
    & \widetilde{M\otimes_S N\otimes_S P}
  }
\tag{2.5.11.4}
\]
